\documentclass[11pt]{scrartcl}
\usepackage[sexy]{evan} %BLBLBLBLBLBLBLBLBLBLBLB EVANCHEN
\usepackage[T1]{fontenc}
\usepackage[utf8]{inputenc}
\usepackage{graphicx}
\usepackage{amsmath}
\usepackage{tikz}
\usepackage{tikz-cd}
\usepackage{asymptote}
\usepackage{amsthm}
\usepackage{amssymb}
\usepackage{gensymb}
\usepackage{amsfonts}
\usepackage[english,italian]{babel}
\usepackage{quoting}
\quotingsetup{font=big}
\usepackage{booktabs}
\usepackage{caption}
\usepackage{lmodern}
\usepackage{faktor}
\usepackage{bbm} %oer fare gli insiemi indicatori
\usepackage{leftindex} %per gli ideali e gli indici a sinistra
\usepackage{hyperref}
\usepackage{epigraph} 
\usepackage{pgfplots}
\usepackage[export]{adjustbox}

%Pacchetti per i codici
\usepackage{listings}
\usepackage{xcolor}
\usepackage{algpseudocode}

%Analisi e probabilità
%----------------------------------
\newcommand{\dif}{\text{d}}
%%% the pure symbol, see https://tex.stackexchange.com/a/718194/4427
\newcommand{\cind}{\perp\mspace{-9mu}\perp}
%%% define \indep
\NewDocumentCommand{\indep}{e{_}}{%
	\mathrel{%
		\IfNoValueTF{#1}{% no subscript
			\cind
		}{% subscript
			\mathchoice{\mathop{\cind}\limits_{#1}}{\cind_{#1}}{\cind_{#1}}{\cind_{#1}}%
		}%
	}%
}
%-----------------------------------

%Algebra lineare e geometria
%----------------------------------
\newcommand{\tr}[1]{\text{tr}\!\left(#1\right)}
\renewcommand{\rank}[1]{\text{rank}\!\left(#1\right)}
\newcommand{\Immagine}[1]{\text{Im}\left(#1\right)}
\renewcommand{\ker}[1]{\text{Ker}\left(#1\right)}
\renewcommand{\AA}{\mathbb{A}}
\newcommand{\LL}{\mathbb{L}}
\newcommand{\SO}{\text{SO}}
%----------------------------------

%Fisica
%----------------------------------
\newcommand{\UDiMisura}[1]{\,\text{#1}}
\newcommand{\ihat}{\hat{\imath}}
\newcommand{\jhat}{\hat{\jmath}}
\newcommand{\khat}{\hat{k}}
\newcommand{\posit}{\vec{r}}
\newcommand{\veloc}{\dot{\vec{r}}}
\newcommand{\accel}{\ddot{\vec{r}}}
%----------------------------------

%Informatica
%----------------------------------
\definecolor{codegreen}{rgb}{0,0.6,0}
\definecolor{codegray}{rgb}{0.5,0.5,0.5}
\definecolor{codepurple}{rgb}{0.58,0,0.82}
\definecolor{backcolour}{rgb}{0.95,0.95,0.92}

\lstdefinestyle{overleafwiki}{
	backgroundcolor=\color{backcolour},   
	commentstyle=\color{codegreen},
	keywordstyle=\color{magenta},
	numberstyle=\tiny\color{codegray},
	stringstyle=\color{codepurple},
	basicstyle=\ttfamily\footnotesize,
	breakatwhitespace=false,         
	breaklines=true,                 
	captionpos=b,                    
	keepspaces=true,                 
	numbers=left,                    
	numbersep=5pt,                  
	showspaces=false,                
	showstringspaces=false,
	showtabs=false,                  
	tabsize=2
}

\lstset{style=overleafwiki}
\newcommand{\bigo}[1]{\mathcal{O}\!\left(#1\right)}
\newcommand{\bigomega}[1]{\Omega\!\left(#1\right)}
\newcommand{\bigtheta}[1]{\Theta\!\left(#1\right)}
%----------------------------------

%Insiemistica e Algebra
%----------------------------------
\newcommand{\eset}{\varnothing}
\newcommand{\parti}[1]{\mathcal{P}\!\left(#1\right)}
\newcommand{\inj}{\hookrightarrow}
\newcommand{\surj}{\twoheadrightarrow}
\newcommand{\bij}{\hookrightarrow\mathrel{\mspace{-15mu}}\rightarrow}
\newcommand{\ordine}[2]{\text{ord}_{#1}\left(#2\right)}
\newcommand{\normal}{\mathrel{\unlhd}}
\newcommand{\MCD}{\text{MCD}}
\renewcommand{\mcm}{\text{mcm}}
%----------------------------------


\begin{document}
	\title{
		Group Representation Theory Notes
		 }
	\subtitle{
		Giovanna Carnovale - SGSS \\
		Andriy Regeta - SGSS 
	}
	\date{Ottobre - Dicembre 2026}
	
	
	\maketitle
	
	\tableofcontents
	
	\newpage
	
	\section{Basics of Representation theory}
	
	\begin{definition}[Group representation]
		A \textit{representation} of a group $G$ is a morphism $\rho:G\to \GL(V)$, where $V$ is a certain vector space. Commonly one refers to $V$ as being the representation of $G$, taking as granted from the context the mosphism $\rho$. 
	\end{definition}
	Analogously a representation of a ring or a field is a morphism from it to an arbitrary $\GL(V)$. As it's usual with actions, if $\rho$ is clear from the context we write
	$$\rho_V(g)v=g\circ v=gv$$
	with a slight abuse o notation.
	\begin{definition}[Group algebra]
		Given a field $k$ and a group $G$ the \textit{group algebra} of $G$ is the set
		$$kG:=\left\{\sum_{g\in G}\lambda_g g:\lambda_g\in k\right\}$$
		equipped with the obvious sum (as a vector space) and the distributive multiplication that satisfies the multiplication of $G$.
	\end{definition}
	We should prove that the one above is actually an algebra, but it is quite easy indeed. Given a representation $\rho$ of $G$ we can extend it linearly to a representation $\tilde{\rho}$ such that
	$$\tilde{\rho}:kG \to \text{End}(V)$$
	and this is a ring representation (we go in End because multiplication in $kG$ is not necessarily invertible).
	\begin{definition}[Algebra]
		Given a field $k$, an \textit{algebra} over it is a $k$-vector space equipped with a bilinear, unitary and associative product.
	\end{definition}
	Just as with groups we define a representation of an algebra as a morphism $\rho:A \to \text{End}(V)$.
	\begin{example}[Trivial representation]
		For any group the map
		$\rho:G\to \GL(V)$
		such that
		$$g\mapsto 1$$
		is the so called \textit{trivial representation} of $G$.
	\end{example}
	\begin{example}[Regular representation]
		For any group $G$, if we take $V=kG$ we get the so called \textit{regular representation} of $G$ as 
		$$h\circ \sum_{g\in G}\lambda_g g \mapsto \sum_{g\in G} h\lambda_g g.$$
	\end{example}
	\begin{example}[Representation of a permutation]
		For $G=S_n$ and $V=k^n$ the map that permutes the canonical basis is a representation, in other words
		$$g\circ e_i := e_{g(i)}$$
		is a representation.
	\end{example}
	\begin{example}[Determinant]
		Given a representation $\rho:G\to \GL(V)$ the map 
		$$\det\circ \rho$$
		is also a representation (in particular is always $1$-dimensional).
	\end{example}
	
	\subsection{Morphisms of representations}
	
	\begin{definition}[Morphisms of representations]
		Given a group $G$ and two representations $\rho_V$ and $\rho_W$ over the same field $k$, a \textit{morphism} between theese is a linear map $T:V\to W$ such that is compatible with both $\rho_V$ and $\rho_W$; that is, we have the following commutative diagram
		\begin{center}
			\begin{tikzcd}
				V\arrow[r,"T"] \arrow[d,"\rho_V(g)",swap]& W\arrow[d,"\rho_W(g)"] \\
				V\arrow[r,"T",swap] & W
			\end{tikzcd}
		\end{center}
		for any $g\in G$. Theese are also called \textit{intertwined operators} or $G$-\textit{equivariant maps}, or $G$-\textit{morphisms}... they have a lot of names. Anyway we say that a morphism of representations is an \textit{isomorphism} if it is invertible.
	\end{definition}
	What we are saying essentially is, being that $\rho_V(g)$ and $\rho_W(g)$ are matrices, $T$ commutes with all such matrices in the way given above. The set of all intertwined operators is oftend denoted as
	$$\text{Hom}(V,W)\ge \text{Hom}_G(V,W)$$
	which is always a non empty vector space (!). \\
	Let's now consider two representations $V$ and $W$, we can investigate what happens to the action of $G$ when we take $V\oplus W$, as we expected
	$$g\circ (v,w)=\left(g\circ v,g\circ w\right)$$
	and no one is surprised. As usual with vector spaces, given two finite dimensional representations $V,W$, if we choose basis $\mathcal{B}_V$, $\mathcal{B}_W$ theese induce a morphism of representations from $V$ and $W$ to
	$$\rho_V:G\to \GL_n(k)$$
	and
	$$\rho_W:G\to \GL_m(k)$$
	in this way $\rho_{V\oplus W}(g)$ has also a nice matrix representation
	$$\rho_{V\oplus W}(g)=
	\left(\begin{array}{c|c}
		\rho_V(g) & 0 \\
		\hline 
		0 & \rho_W(g)
	\end{array}\right).$$
	\begin{lemma}
		Given a map between two representations $V$ and $W$ of $G$, we have that this is $G$-equivariant iff is $kG$ equivariant, given $k$ the field over which $V$ is defined.
	\end{lemma} 
	\begin{example}[Dual representation]
		Given a representation $V$ of $G$, it's dual space $V^*$ has naturally a representation structure
		$$\left(\rho_{V^*}(g)f\right)(v)=f\left(\rho_V(g)^{-1}v\right)$$
		where the inverse is taken to make sure that it is an homomorphism.
		\begin{center}
			\begin{tikzcd}
				contenuto...
			\end{tikzcd}
		\end{center}
		a more compact way to write it is $g\circ f = f\circ g^{-1}$. An even more compact way to see is through matrices, choosing a basis $B$ and it's dual basis $B^*$ of $V^*$ we see that the actions, as matrices satisfy
		$$\rho_{V^*}(g)=\rho_V(g)^\intercal$$
	\end{example}
	\begin{definition}[Tensor product of representations]
		Given two representations $V$ and $W$ of a group their \textit{tensor product} is the space 
		$$V\otimes W := \langle \left\{v \otimes w: v\in V,w\in W\right\}\rangle$$
		equipped with the action 
		$$\rho_{V\otimes W}(g)(v\otimes w):=\rho_V(g)v\otimes \rho_W(g)w.$$
	\end{definition}
	\begin{remark}
		Given spaces $W$ and $V^*$ there is always a canonical injection
		\begin{align*}
			F: W\otimes V^* & \inj \text{Hom}(V,W) \\
			\sum c_{wf}w\otimes f & \mapsto \sum c_{wf} f(v)w
		\end{align*}
		which is equivalent to say that on pure tensors our map is just 
		$$w\otimes f \mapsto F(w\otimes f)(v):=f(v)w$$
		and for all the other elements of $V\otimes W$ we extend it linearly.  
	\end{remark}
	This gives us immediately that if $V$ and $W$ representations of $G$, so is $V^*\otimes W$. Since the canonical injection defined above is equivariant (prove it) we hence have that $\text{Hom}(V,W)$ is a representation of $G$. What is the action of this? We have that
	$$\left(g\circ f\right)(v)=\rho_W(g)\left(f\left(\rho_V(g^{-1})v\right)\right)$$
	and so if $V=W$ this is just a conjugation of $f\in \text{End}(V)$!
	
	\subsection{Subrepresentations}
	
	\begin{definition}[Subrepresentation]
		Given a representation $V$ of an object $G$, a \textit{subrepresentation} of $V$ is a subspace $U\le V$ such that is $G$-stable, that is 
		$$G\circ U \subseteq U$$
		which means that acting with $G$ on $U$ always gets us to elements still in $U$.
	\end{definition}
	\begin{example}[Fixed point of permutations]
		Given $G=S_n$ with it's permutation representation on $V=k^n$, a subrepresentation of $V$ is
		$$U=\langle e_1+e_2+\dots+e_n\rangle$$
		and in this case is much more than $G$ invariant: is pointwise fixed by the action of $G$. A non pointwise fixed example is the space $W$ with cartesian equation
		$$\sum_{i=1}^{n}x_i=0$$
		which, if $\text{car}K \not \mid n$ is not pointwise fixed for every point, $w\not\in G\circ w$ for any $w$. 
	\end{example}
	An operation we haven't yet touched is the quotient of spaces, we want to investigate what happens if we take two representation $V\ge U$ of $G$ and consider $\faktor{V}{U}$. We claim that this is a representation with the action
	$$\rho_{V/U}(v+U):=\rho_V(g)v+U$$
	and we should check that this is actually well defined and is a homomorphism. What does this mean in matrix terms? Given a basis $u_i$ of $U$ and a completion of it to a basis of $V$, the matrix of $\rho_V(g)$ in this basis has generally the form
	$$\rho_V(g)=
	\left(
	\begin{array}{c|c}
		\rho_U(g) & * \\
		\hline
		0 & \rho_{V/U}(g)
	\end{array}
	\right)
	$$
	one of the big questions of representation theory, as always when we deal with matrices, is wether, given a certain representation $V$ of a group $G$, we can express the action of $G$ on $V$ as acting through block diagonal matrices. From the formula above we see that if we can find a couple of spaces $U\oplus U'$ that are both $G$ stable, we have achieved our goal. 
	\begin{example}[Decomposition of a representation of $V$]
		Taking $G=S_n$ and $V=k^n$, since we have 
		$$\langle e_1+e_2+\dots+e_n \rangle \oplus \left\langle \left\{\sum_{i=1}^n x_i=0:x\in k^n\right\}\right\rangle$$
		and theese are $G$-stable we have that this is a decomposition of $V$ as a representation.
	\end{example}
	
\end{document}
